\documentclass[11pt]{article}

\usepackage[margin=1.05in]{geometry}
\usepackage{lmodern}
\usepackage{amsmath,amssymb}
\usepackage{booktabs}
\usepackage{tabularx}
\usepackage{enumitem}
\usepackage{xcolor}
\usepackage{graphicx}
\usepackage{caption}
\usepackage[expansion=false]{microtype}
\usepackage[colorlinks=true,linkcolor=blue!55!black,citecolor=blue!50!black,urlcolor=blue!55!black]{hyperref}

\newcommand{\code}[1]{\texttt{#1}}
\newcommand{\relL}{\text{rel-}L_2}
\newcolumntype{Y}{>{\raggedright\arraybackslash}X}

\title{\bfseries Calibrating Self-Exciting Event Counts with a Learned Forward
Operator\\[3pt]\large Technical report}
\author{Alberto Gennaro}
\date{24 June 2026}

\begin{document}
\maketitle

\section{Scope}

This report describes a method for fitting multivariate self-exciting (Hawkes) models
to \emph{interval-censored} event data --- counts aggregated into fixed time bins ---
together with bin-level covariates, and the role a learned neural operator plays in
that fit. Section~\ref{sec:problem} states the data and the model. Section~\ref{sec:pipeline}
is the core of the report: it gives the calibration procedure step by step for a
concrete, common case (covariates in the baseline only, constant excitation, weekly
bins). Section~\ref{sec:operator} defines the forward operator and reports its accuracy.
Sections~\ref{sec:inverse}--\ref{sec:application} cover parameter recovery and the
downstream ranking application. Section~\ref{sec:limits} lists limitations.
All results are on synthetic data at CPU scale.

\section{Problem setting and model}
\label{sec:problem}

\paragraph{Data.} There are $M$ event streams (``sectors''), observed over $W$ time bins
(e.g.\ weeks) with endpoints $0=a_0<a_1<\dots<a_W=T$. For each bin $i$ and sector $m$
we observe an event count $C_{i,m}\in\mathbb{N}$. We also observe a covariate vector
$X_i\in\mathbb{R}^p$ per bin (for example a popularity score, a sector indicator, a
macro regime). Event times within a bin are \emph{not} observed.

\paragraph{Model.} Each sector $m$ is driven by a conditional intensity
\begin{equation}
\lambda_m(t) = s_m(t) + \sum_{j=1}^{M} \sum_{t_k^{(j)} < t} A_{m,j}\, e^{-\theta\,(t-t_k^{(j)})},
\label{eq:hawkes}
\end{equation}
with a covariate-driven, log-linear baseline and a constant excitation matrix:
\begin{equation}
s_m(t) = \exp\!\big(\gamma_{m,0} + \gamma_m^\top X(t)\big),
\qquad X(t)=X_i \ \text{for } t\in(a_{i-1},a_i],
\label{eq:baseline}
\end{equation}
and $A_{m,j}\ge 0$ the excitation of sector $m$ by sector $j$, decaying at rate
$\theta$. In this report the excitation $A$ does \emph{not} depend on covariates; the
covariates enter only through the baseline. (The covariate-modulated-excitation case is
a direct extension, noted in Section~\ref{sec:limits}.)

\paragraph{The quantity we can compute from counts.} The Hawkes likelihood needs event
times, which we do not have. We therefore work with the Mean Behavior Poisson Process
(MBPP): the expected intensity $\xi=\mathbb{E}[\lambda]$ is deterministic and solves the
Volterra equation
\begin{equation}
\xi(t) = s(t) + \int_0^t \Phi(t-u)\,\xi(u)\,\mathrm{d}u,
\qquad \Phi_{m,j}(\tau)=A_{m,j}\,e^{-\theta\tau}.
\label{eq:mbpp}
\end{equation}
By Watanabe's characterization the process with deterministic compensator
$\Xi(t)=\int_0^t\xi$ is an inhomogeneous Poisson process, so the bin counts are
independent Poisson variables,
\begin{equation}
C_{i,m}\sim \mathrm{Poisson}\big(\Xi_{i,m}\big),\qquad
\Xi_{i,m}=\int_{a_{i-1}}^{a_i}\xi_m(t)\,\mathrm{d}t .
\label{eq:poisson}
\end{equation}
Calibration is maximum likelihood for this Poisson model. Up to a constant, the
negative log-likelihood in the parameters $\Theta=(\gamma,A,\theta)$ is
\begin{equation}
\mathcal{L}(\Theta) = \sum_{i=1}^{W}\sum_{m=1}^{M}
\Big(\Xi_{i,m}(\Theta) - C_{i,m}\,\log \Xi_{i,m}(\Theta)\Big).
\label{eq:nll}
\end{equation}

\section{Calibration procedure}
\label{sec:pipeline}

Evaluating $\mathcal{L}$ in \eqref{eq:nll} requires the bin compensators
$\Xi_{i,m}(\Theta)$, which come from solving \eqref{eq:mbpp} for the current
parameters. That forward solve is the only expensive step, and it is repeated at every
optimiser iteration. The method replaces it with a single call to a forward operator
that has been trained, once, to solve \eqref{eq:mbpp} for any instance in the relevant
family.

\paragraph{The forward operator.} Let $G_\varphi$ denote a neural operator (weights
$\varphi$) that approximates the solution map
\begin{equation}
G_\varphi:\ \big(s(\cdot),\,A\big)\ \longmapsto\ \xi(\cdot),
\label{eq:operator}
\end{equation}
i.e.\ it maps a baseline path $s$ (sampled on a fixed time grid) and an excitation
matrix $A$ to the mean-intensity path $\xi$ on that grid, in one forward pass.
Section~\ref{sec:operator} describes how $G_\varphi$ is built and trained; here we use
it as a black-box, differentiable solver.

\paragraph{Offline (once).} Train $G_\varphi$ over the distribution of instances the fit
will query: baseline paths of the form \eqref{eq:baseline} spanning the plausible range
of $(\gamma,X)$, and subcritical excitation matrices $A$. Training minimises the
residual of \eqref{eq:mbpp} (Section~\ref{sec:operator}); the trained operator is then
reused for every dataset.

\paragraph{Online (per dataset).} Fit $\Theta$ by gradient descent on
\eqref{eq:nll}, using $G_\varphi$ as the forward map. One iteration:
\begin{enumerate}[leftmargin=1.7em,itemsep=2pt]
\item \textbf{Baseline path from covariates.} With the current $\gamma$, evaluate the
piecewise-constant baseline on the grid, $s_m(t_k)=\exp(\gamma_{m,0}+\gamma_m^\top
X(t_k))$, using the observed weekly covariates $X_i$.
\item \textbf{Forward solve.} $\xi = G_\varphi(s, A)$ --- a single network pass, no
integral-equation solve.
\item \textbf{Bin the compensator.} $\Xi_{i,m}=\int_{a_{i-1}}^{a_i}\xi_m$, computed by
the trapezoid rule on the grid.
\item \textbf{Likelihood.} Evaluate $\mathcal{L}(\Theta)$ from \eqref{eq:nll} against the
observed counts $C_{i,m}$.
\item \textbf{Gradient and step.} Because $G_\varphi$ is differentiable and $s$ depends
on $\gamma$ in closed form, obtain $\nabla_\Theta\mathcal{L}$ by automatic
differentiation through the whole chain (no finite differences, no solver loop) and take
a descent step.
\end{enumerate}
The output is the fitted baseline coefficients $\gamma$ (the covariate effects), the
excitation matrix $A$, and the decay $\theta$. Standard errors follow from the observed
information; a Bayesian posterior is available by sampling \eqref{eq:nll} plus a prior.

\paragraph{What the operator does and does not change.} The operator changes only how the
forward solve in step~2 is computed; the statistical model \eqref{eq:hawkes}--\eqref{eq:nll}
and the estimator are unchanged. For the baseline-only case of this section the forward
solve also has an exact, inexpensive numerical solution (the kernel is a convolution),
so the operator is used here for speed and for clean autodiff gradients; its decisive
advantage appears in the covariate-modulated-excitation case (Section~\ref{sec:limits}),
where the equation is no longer a convolution and the per-iteration solve is the
bottleneck.

\section{The forward operator}
\label{sec:operator}

\paragraph{Architecture.} $G_\varphi$ is a DeepONet: a branch network encodes the
instance --- the baseline path $s$ sampled on the grid and the excitation matrix $A$ ---
and a trunk network encodes the query time; their inner product gives $\xi(t)$. The
operator must be conditioned on the instance it is solving: a network that does not
receive $A$ (or receives only the forcing while the system varies) is not a function on
the family --- the same input corresponds to many correct outputs --- and cannot be
fit below roughly $35$--$50\%$ relative error. Conditioning on $(s,A)$ removes this.

\paragraph{Training objective.} On a grid $\{t_k\}$, write the discretised version of
\eqref{eq:mbpp} as $\xi = s + A\,(G\xi)$, where $G$ is the lower-triangular matrix
implementing $\int_0^t e^{-\theta(t-u)}(\cdot)\,\mathrm{d}u$. The residual
\begin{equation}
R[\xi] = \xi - s - A\,(G\xi)
\label{eq:residual}
\end{equation}
is linear in $\xi$, so it has a unique zero at the exact solution. The operator is
trained to drive $R$ to zero,
\begin{equation}
\min_\varphi\ \mathbb{E}_{(s,A)}\!\left[\frac{\lVert R[G_\varphi(s,A)]\rVert^2}{\lVert
G_\varphi(s,A)\rVert^2}\right] \;+\; (\text{supervised term on a set of exact anchor
solutions}),
\label{eq:trainloss}
\end{equation}
where the relative normalisation keeps large-intensity instances from dominating, and
the anchor term (a small set of solutions from the exact solver) accelerates and
stabilises training. No labelled solutions are required for the residual term.

\paragraph{Accuracy.} Trained over instances of dimension $M$ and evaluated against the
exact multivariate solver on $400$ held-out instances per dimension:

\begin{table}[h]
\centering\small
\begin{tabular}{@{}l r r r r r@{}}
\toprule
$M$ (sectors) & instance dim ($M{+}M^2$) & $\relL$ mean & median & p90 & training time \\
\midrule
3 & 12 & 0.55\% & 0.47\% & 0.87\% & 118 s \\
5 & 30 & 1.52\% & 1.36\% & 2.40\% & 161 s \\
8 & 72 & 2.03\% & 1.85\% & 3.11\% & 242 s \\
\bottomrule
\end{tabular}
\caption{Relative $L_2$ error of $G_\varphi$ against the exact solver, held out. Error is
approximately constant across the branching ratio and uniform across sectors. One
evaluation takes about $3$~ms. Source: \code{experiments/exp17\_pino\_jax.py}.}
\label{tab:operator}
\end{table}

Trained on noise-corrupted targets and tested against the clean solver, the error
remains within twice its noiseless value up to about $25\%$ multiplicative noise.

\begin{figure}[h]
\centering
\includegraphics[width=0.9\linewidth]{../results/exp17_pino_jax_summary.png}
\caption{Forward operator versus the exact multivariate solver, by dimension and
branching ratio.}
\label{fig:operator}
\end{figure}

\section{Parameter recovery}
\label{sec:inverse}

Two routes recover the parameters from an observed intensity (or, through
\eqref{eq:poisson}, from counts).

The first is the likelihood fit of Section~\ref{sec:pipeline} itself, which returns
$(\gamma,A,\theta)$ with standard errors. The second is direct inversion: minimise
$\lVert G_\varphi(s,A)-\xi_{\mathrm{obs}}\rVert^2/\lVert\xi_{\mathrm{obs}}\rVert^2$ over
the parameters by gradient descent through the operator. On clean data both the
covariate coefficients and the excitation matrix are recovered to within a few percent
(covariate coefficients to $\le 0.1\%$, excitation to $\le 2.5\%$ at $M\le 5$). Under
observation noise the covariate coefficients remain the well-determined part (within
$5\%$ at noise level $\sigma=0.02$ on a single path), while the excitation matrix is
more sensitive and is recovered by averaging $R$ repeated observations, with error
scaling as $\sigma/\sqrt{R}$. A one-shot regression network that maps counts directly
to parameters is fast but only weakly accurate (entry correlation about $0.38$, overall
stability correlation about $0.48$) and is suitable as an initialisation, not as a final
estimate. Source: \code{experiments/exp18\_covariate\_inverse.py}.

\section{Application: ranking the next event}
\label{sec:application}

The downstream task is to rank the sectors (or firms) most likely to produce an event
within a horizon $h$. Given a fitted model, the rank score is the probability of at
least one event of stream $m$ in $(t,t+h]$,
\begin{equation}
p_m(t,h) = 1 - \exp\!\Big(-\!\int_t^{t+h}\xi_m(u)\,\mathrm{d}u\Big),
\label{eq:rank}
\end{equation}
the compensator increment, which the forward operator already produces.

An alternative model class for the same task is survival analysis with a ranking
objective. The time to the next event is modelled by a discrete-time competing-risks
hazard; the model is fitted by the survival likelihood (a masked Bernoulli over time
bins) plus a pairwise ranking loss on the cumulative-incidence function, and the
population baseline is the Kaplan--Meier curve. On a synthetic firm population the
fitted model attains a concordance index of $0.687$ against $0.617$ for a
single-covariate baseline ($0.5$ is random), with calibrated survival curves and a
Recall@$k$ improvement over the baseline (Figure~\ref{fig:survival}). The relationship
to the Hawkes model is that a cause-specific hazard is a per-stream intensity; survival
lets it depend on covariates but not on the event history of other streams, which the
Hawkes excitation adds. Detail: \code{SURVIVAL\_RANKING.md},
\code{experiments/exp16\_survival.py}.

\begin{figure}[h]
\centering
\includegraphics[width=0.9\linewidth]{../results/exp16_survival.png}
\caption{Survival model with a ranking objective: fitted survival curves against the
Kaplan--Meier baseline (a); concordance index (b); predicted risk by outcome (c);
Recall@$k$ (d).}
\label{fig:survival}
\end{figure}

\section{Software and reproducibility}

The model, fitters, exact solvers, operator, and experiments are in the
\code{hawkes\_calibration} package. The counts-based fitters are
\code{fit\_mbpp\_ic\_covariates} (baseline covariates) and
\code{fit\_mbpp\_ic\_excitation\_multi} (covariate-modulated excitation, multivariate);
the forward operator is in \code{hawkes\_calibration/operators}; goodness-of-fit
(time-rescaling) and a Bayesian posterior are provided. The automated test suite passes
($53/53$). Reproduction commands accompany each experiment script.

\section{Limitations}
\label{sec:limits}

\begin{itemize}[leftmargin=1.3em,itemsep=2pt]
\item Results are on synthetic data at CPU scale (small networks, thousands of training
instances). Validation at larger scale and on real binned-count series is required
before drawing conclusions about absolute accuracy.
\item Covariate-modulated excitation, $A_{m,j}(t)=A_{m,j}\exp(\delta^\top Z(t))$, uses
the same procedure as Section~\ref{sec:pipeline} with the operator additionally
conditioned on $(\delta, Z)$; this is the case where the per-iteration solve is most
costly and the operator most useful. It is implemented in the exact solver and fitter;
the operator for this case is the immediate next step.
\item The branching ratio (overall stability) is well determined from counts; the decay
$\theta$ and the individual cross-excitation entries are weakly identified, the more so
the coarser the bins.
\item For $M$ beyond roughly $15$, the $M^2$ excitation matrix should be encoded
through a low-rank parameterisation rather than passed entry-wise to the operator.
\end{itemize}

\section{References}

Checked against source (authors, venue, year, link) on 24 June 2026.

\subsection*{Model and estimation}
\begin{enumerate}[leftmargin=1.6em,itemsep=2pt]
\item Rizoiu, Soen, Li, Calderon, Dong, Menon, Xie (2022). Interval-censored Hawkes
processes. JMLR 23(1):1--84. \url{https://www.jmlr.org/papers/v23/21-0917.html}
\item Hawkes (1971). Spectra of some self-exciting and mutually exciting point
processes. Biometrika 58(1):83--90.
\item Hawkes \& Oakes (1974). A cluster process representation of a self-exciting
process. J.\ Appl.\ Probab.\ 11(3).
\item Daley \& Vere-Jones (2003). An Introduction to the Theory of Point Processes,
Vol.~I. Springer.
\item Banerjee, Merugu, Dhillon, Ghosh (2005). Clustering with Bregman divergences. JMLR 6.
\end{enumerate}

\subsection*{Neural operators and point processes}
\begin{enumerate}[leftmargin=1.6em,itemsep=2pt,start=6]
\item Li, Kovachki, Azizzadenesheli, Liu, Bhattacharya, Stuart, Anandkumar (2021).
Physics-Informed Neural Operator for Learning PDEs. \url{https://arxiv.org/abs/2111.03794}
\item Lu, Jin, Pang, Zhang, Karniadakis (2021). Learning nonlinear operators via DeepONet.
Nature Machine Intelligence 3:218--229.
\item Du, Dai, Trivedi, Upadhyay, Gomez-Rodriguez, Song (2016). Recurrent Marked Temporal
Point Processes. KDD. \url{https://www.kdd.org/kdd2016/papers/files/rpp1081-duA.pdf}
\item Mei \& Eisner (2017). The Neural Hawkes Process. NeurIPS. \url{https://arxiv.org/abs/1612.09328}
\item Zuo, Jiang, Li, Zhao, Zha (2020). Transformer Hawkes Process. ICML (PMLR 119).
\url{http://proceedings.mlr.press/v119/zuo20a/zuo20a.pdf}
\item Zhang, Lipani, Kirnap, Yilmaz (2020). Self-Attentive Hawkes Processes. ICML (PMLR 119).
\url{https://arxiv.org/abs/1907.07561}
\item Shchur, Bilo\v{s}, G\"unnemann (2020). Intensity-Free Learning of Temporal Point
Processes. ICLR. \url{https://arxiv.org/abs/1909.12127}
\item Trivedi, Dai, Wang, Song (2017). Know-Evolve: Deep Temporal Reasoning for Dynamic
Knowledge Graphs. ICML (PMLR 70). \url{https://arxiv.org/abs/1705.05742}
\item Trivedi, Farajtabar, Biswal, Zha (2019). DyRep: Learning Representations over
Dynamic Graphs. ICLR. \url{https://openreview.net/forum?id=HyePrhR5KX}
\item Xue et al.\ (2024). EasyTPP: Towards Open Benchmarking Temporal Point Processes.
ICLR. \url{https://arxiv.org/abs/2307.08097}
\end{enumerate}

\subsection*{Survival analysis and ranking}
\begin{enumerate}[leftmargin=1.6em,itemsep=2pt,start=16]
\item Kaplan \& Meier (1958). Nonparametric estimation from incomplete observations. JASA
53:457--481.
\item Cox (1972). Regression models and life-tables. JRSS-B 34:187--220.
\item Fine \& Gray (1999). A proportional hazards model for the subdistribution of a
competing risk. JASA 94:496--509.
\item Graf, Schmoor, Sauerbrei, Schumacher (1999). Assessment and comparison of prognostic
classification schemes for survival data. Statistics in Medicine 18(17--18):2529--2545.
\item Antolini, Boracchi, Biganzoli (2005). A time-dependent discrimination index for
survival data. Statistics in Medicine 24(24):3927--3944.
\url{https://onlinelibrary.wiley.com/doi/10.1002/sim.2427}
\item Katzman, Shaham, Cloninger, Bates, Jiang, Kluger (2018). DeepSurv. BMC Medical
Research Methodology 18:24. \url{https://link.springer.com/article/10.1186/s12874-018-0482-1}
\item Lee, Zame, Yoon, van der Schaar (2018). DeepHit. AAAI.
\url{https://cdn.aaai.org/ojs/11842/11842-13-15370-1-2-20201228.pdf}
\item Lee, Yoon, van der Schaar (2019). Dynamic-DeepHit. IEEE Trans.\ Biomedical
Engineering 67(1).
\item Ross, Das, Sciro, Raza (2021). CapitalVX: a machine learning model for startup
selection and exit prediction. J.\ Finance and Data Science.
\url{https://www.sciencedirect.com/science/article/pii/S2405918821000040}
\end{enumerate}

\paragraph{Reference implementations.} \code{pycox} and \code{scikit-survival} (DeepHit,
Cox, Brier, concordance); \code{EasyTPP} (RMTPP / NHP / THP baselines).

\section*{Related documents}
\code{paper/investment\_case\_study.pdf} (the model and proofs);
\code{SURVIVAL\_RANKING.md} (survival method in full);
\code{NEXT\_FUNDING\_DESIGN.md} (the ranking application);
\code{README.md} (the package).

\end{document}
